\chapter{Локальное причинное ядро распространения ванны}
\label{ch:v10-nonequilibrium-bath-local-causal-kernel}

После отрицательного аудита готовых спектральных профилей естественно
выводить память из локального правила распространения. Пусть $A$ ---
симметричный оператор ближайших соседей клеточного графа,
$A_{ij}=1/2$ при $d(i,j)=1$ и ноль иначе. Тогда
\begin{equation}
 (A^n)_{ij}=0\qquad\text{при}\qquad d(i,j)>n.
\end{equation}
Это точный дискретный световой конус: за один шаг информация проходит не
более одного ребра.

Минимальное затухающее семейство имеет вид
\begin{equation}
 K_n(r)=r^n A^n,\qquad 0<r<1.
\end{equation}
Его огибающая геометрична, а полное число шагов памяти равно
\begin{equation}
 M(r)=\sum_{n=0}^{\infty}r^n=\frac{1}{1-r}.
\end{equation}
Условное сопоставление двух прежних KMS-отношений даёт
\begin{equation}
 r_h=\frac12,\quad M_h=2,
 \qquad
 r_c=\frac14,\quad M_c=\frac43.
\end{equation}
Трёхточечные ковариационные матрицы
$C_{ij}=r^{|i-j|}$ имеют определители $9/16$ и $225/256$, поэтому оба
свидетеля положительны.

Но локальность не выбирает $r$. Цепной родитель
\begin{equation}
 \mathcal P_K=\frac12\bigl[(k_1-r)^2+(k_2-rk_1)^2+(k_3-rk_2)^2\bigr]
\end{equation}
обнуляется на всём семействе $(k_1,k_2,k_3)=(r,r^2,r^3)$. При $r=1/2$
его гессиан имеет ранг/ядро $3/1$, а вектор ядра
$(1,1,3/4,1)$ является касательной к этой кривой. Поэтому причинность
фиксирует поддержку ядра, но не его затухание.

Для переменных
$(\tau_{\rm corr},\Delta t,\ell_{\rm cell},v_g)$ размерная карта имеет
ранг/ядро $2/2$. Фиксация длительности шага оставляет одну орбиту; только
дополнительный независимый якорь длины закрывает карту.

\begin{theorem}[Локальная причинная семья памяти]
Оператор ближайших соседей порождает ядра, строго поддержанные внутри
клеточного светового конуса, а положительная геометрическая огибающая
$r^n$ даёт конечную память $1/(1-r)$ для каждого $0<r<1$.
\end{theorem}

\begin{nogocriterion}[Причинность не является диссипацией]
Все значения $0<r<1$ имеют одинаковую локальную поддержку. Отождествление
$r$ с отношением KMS-переходов допустимо только после отдельного
типизированного морфизма; без него коэффициент затухания и абсолютное время
памяти не выведены.
\end{nogocriterion}

\begin{protocol}\begin{enumerate}
\item условная локально-причинная архитектура: \textbf{$8/8$};
\item проверенные шаги причинной поддержки: \textbf{$3/3$};
\item происхождение коэффициента затухания: \textbf{$0/1$};
\item происхождение абсолютного времени: \textbf{$0/1$};
\item следующий гейт: \textbf{типизированное вложение ядра в клеточный комплекс}.
\end{enumerate}\end{protocol}

Машинный сертификат сохранён в
\path{s2t_v10_cell_birth_four_volume_nonequilibrium_bath_local_causal_propagation_kernel_parent_origin_gate_results.json}.